\chapter{Physical Model Requirements}

\section{Why PSCs Need Mechanistic Models}

Stability assessment for PSCs is difficult because degradation depends on
device architecture, stress protocol, preconditioning, and measurement history.
The ISOS consensus work formalized how stability tests should be reported, but
standardized testing does not remove the need for mechanism-aware models
\cite{Khenkin_2020}. A useful model must connect an observed electrical change
to a plausible internal cause without hiding uncertainty.

Drift-diffusion modeling is the natural starting point because it represents
carrier generation, recombination, transport, and electrostatics in the device
stack. Recent PSC modeling reviews emphasize that these models remain valuable
when assumptions are explicit and when electronic and ionic degrees of freedom
are not confused \cite{Singh_2025}. COMSOL is used here as the finite-element
environment for that device model, not as a black box.

\section{p-i-n Device Operation and J--V Observables}

The laboratory device is an inverted p-i-n architecture. Light enters through
the glass/ITO side, followed by the p-selective NiO$_x$/MeO-2PACz contact,
the nominally intrinsic metal-halide perovskite absorber, the n-selective
C$_{60}$/BCP contact, and the Ag electrode. Photons absorbed in the
perovskite generate electron--hole pairs. Holes are extracted toward
NiO$_x$/MeO-2PACz and ITO, while electrons are extracted toward C$_{60}$/BCP
and Ag.

The measured J--V curve compresses several physical processes into four common
metrics. The short-circuit current density $J_{sc}$ is most sensitive to
absorption and carrier collection. The open-circuit voltage $V_{oc}$ is
strongly constrained by recombination and energetic alignment. Fill factor
responds to recombination, transport barriers, series resistance, and shunt
leakage. PCE combines all three:
\begin{equation}
\mathrm{PCE}=\frac{V_{oc}J_{sc}\mathrm{FF}}{P_{\mathrm{in}}}.
\end{equation}
Because different internal changes can produce similar metric changes, the
full curve shape, scan direction, scan rate, and measurement history are
required for mechanism-aware interpretation.

\section{Electronic Transport Equations}

The electronic device physics are governed by Poisson electrostatics, carrier
continuity, and electron and hole current-density relations. The electric
potential $\phi$ defines the electric field
\begin{equation}
\mathbf{E}=-\nabla \phi .
\end{equation}
The electron and hole current densities are written as
\begin{align}
\mathbf{J}_n &= q\mu_n n\mathbf{E} + qD_n\nabla n,\\
\mathbf{J}_p &= q\mu_p p\mathbf{E} - qD_p\nabla p,
\end{align}
where $q$ is the elementary charge, $n$ and $p$ are electron and hole
concentrations, $\mu_n$ and $\mu_p$ are mobilities, and $D_n$ and $D_p$ are
diffusion coefficients. The drift terms follow the force on mobile charges
under $\mathbf{E}$, while the diffusion terms follow concentration gradients.
The opposite signs in the diffusion terms reflect the charge sign convention
for electrons and holes.

Carrier conservation is represented by
\begin{align}
\frac{\partial n}{\partial t}
&= \frac{1}{q}\nabla\cdot \mathbf{J}_n + G - R,\\
\frac{\partial p}{\partial t}
&= -\frac{1}{q}\nabla\cdot \mathbf{J}_p + G - R,
\end{align}
where $G$ is photogeneration and $R$ is the total recombination rate. The
different signs in the divergence terms follow from the definitions of
electron and hole conventional current.

Electrostatics is coupled through Poisson's equation:
\begin{equation}
\nabla\cdot(\epsilon\nabla\phi)
= -q\left(p-n+N_D^+-N_A^-+\frac{\rho_{\mathrm{ion}}}{q}\right).
\label{eq:poisson}
\end{equation}
Here $\epsilon$ is the permittivity, $N_D^+$ and $N_A^-$ are ionized donor and
acceptor densities, and $\rho_{\mathrm{ion}}$ is the mobile-ion space-charge
density in C/m$^3$. The expression $\rho_{\mathrm{ion}}/q$ converts ionic
charge density into an equivalent signed number density before multiplication
by $q$.

\section{Generation and Recombination Terms}

Illumination enters the Semiconductor Module as a volumetric electron-hole pair
generation rate $G$ in the perovskite absorber. The optical field is not solved
with a Wave Optics interface. Instead, the model uses an algebraic
Beer--Lambert generation profile that supplies a spatially varying source term
to the semiconductor equations. The photon flux entering the absorber is
represented as
\begin{equation}
\Phi(x)=\Phi_0 e^{-\alpha x},
\end{equation}
where $\Phi_0$ is the incident photon flux at the illuminated side of the
absorber, $\alpha$ is an effective absorption coefficient, and $x$ is the local
depth coordinate in the perovskite. The corresponding generation profile is
proportional to the decaying photon flux and is normalized so that the
absorber-averaged generation remains controlled. This approach preserves a
depth-dependent generation shape without requiring a full optical transfer
matrix. It does not resolve wavelength-dependent absorption, parasitic
absorption in transport layers, interference effects, or optical reflection.

The total electronic recombination rate can be organized as
\begin{equation}
R = R_{\mathrm{rad}}+R_{\mathrm{SRH}}+R_{\mathrm{Auger}}.
\end{equation}
Direct or bimolecular radiative recombination is represented by
\begin{equation}
R_{\mathrm{rad}}=B\left(np-n_i^2\right),
\end{equation}
where $B$ is the radiative recombination coefficient. Auger recombination can
be written as
\begin{equation}
R_{\mathrm{Auger}}=
\left(C_n n+C_p p\right)\left(np-n_i^2\right),
\end{equation}
where $C_n$ and $C_p$ are electron- and hole-assisted Auger coefficients.
Radiative and Auger terms define the complete recombination hierarchy presented
for the device physics, but the current degradation implementation is
calibrated primarily through SRH recombination. Future quantitative use of
$B$, $C_n$, and $C_p$ requires architecture-specific transient or
photoluminescence calibration.

Trap-assisted recombination is represented with the Shockley--Read--Hall (SRH)
form
\begin{equation}
R_{\mathrm{SRH}} =
\frac{np-n_i^2}
{\tau_p(n+n_1)+\tau_n(p+p_1)}.
\label{eq:srh}
\end{equation}
The carrier lifetimes depend on trap density:
\begin{align}
\tau_n &= \frac{1}{\sigma_n v_{\mathrm{th}} N_t},\\
\tau_p &= \frac{1}{\sigma_p v_{\mathrm{th}} N_t},
\end{align}
where $\sigma_n$ and $\sigma_p$ are capture cross sections, $v_{\mathrm{th}}$
is thermal velocity, and $N_t$ is trap density. Increasing $N_t$ reduces
carrier lifetime, increases nonradiative recombination, lowers $V_{oc}$, and
can reduce fill factor or $J_{sc}$ when recombination becomes strong.

\section{Hysteresis as a Modeling Requirement}

J--V hysteresis is an observable result of coupling among mobile ions, internal
electric fields, recombination, and voltage protocol. Huang et al. showed that
dielectric properties and scan rate can strongly affect PSC hysteresis, and
their time-dependent ion-migration model links those effects to built-in field
changes \cite{Huang_2022}. Wang et al. further argued that hysteresis can be
used diagnostically because evolving ion-induced fields interact with carrier
transport and recombination \cite{WangNatComm_2024}.

A model that only produces one steady J--V curve at one scan rate cannot test
whether the internal state is physically credible. Slower scans allow ionic
states to respond more fully; faster scans probe a more frozen internal state.
Agreement across scan protocols is therefore a stronger model-quality test than
agreement at one operating point.

\section{Model Requirements}

The COMSOL model uses the following requirements:
\begin{enumerate}
\item The device stack and contacts must be explicit enough for comparison to
laboratory samples.
\item Voltage protocol variables must be treated as controlled inputs, not as
hidden physics.
\item Internal material and degradation states must be represented as model
parameters or state variables, not inferred from PCE alone.
\item Mobile-ion redistribution must enter Poisson electrostatics through a
charge-density term with consistent units.
\item Degradation variables must map to COMSOL-accessible quantities such as
$N_t$, $R_s$, $R_{sh}$, surface recombination, and photogeneration.
\item Model claims must be limited to the parameter range, stress protocol, and
calibration state actually simulated.
\end{enumerate}
